% =============================================================================
% HTGR Simulator Demonstration
% =============================================================================

\documentclass[aspectratio=169]{beamer}

\usepackage{beamerthememidcenturymodern}

% Theme
\mcmTheme{Kraft}

\institute{Singapore Nuclear Research and Safety Initiative}

\title[HTGR Simulator]{Learn About High Temperature Gas-Cooled Reactors}
\subtitle{A Live Educational Simulation}
\author{Dr Theodore Ong}
\date{8 September 2026}

\begin{document}

% =============================================================================
\begin{frame}
    \titlepage
\end{frame}

% =============================================================================
\begin{frame}{What is a High Temperature Gas-Cooled Reactor (HTGR)?}

\begin{columns}

\begin{column}{0.55\textwidth}

\begin{itemize}
    \item Uses helium gas as coolant
    \item Uses graphite to moderate neutrons
    \item Produces heat through nuclear fission
    \item Heat is converted into electricity
    \item Designed with strong passive safety features
\end{itemize}

\vspace{1em}

\begin{block}{Key Idea}
A nuclear reactor is simply a large heat source.

The heat can be used to generate electricity.
\end{block}

\end{column}

\begin{column}{0.45\textwidth}

\centering

% Replace this with your schematic screenshot
\includegraphics[width=\linewidth]{htr_sim_full.png}

\end{column}

\end{columns}

\end{frame}

% =============================================================================
\begin{frame}{How Does Heat Become Electricity?}

\centering

\Huge

Reactor

\vspace{0.4em}

$\Downarrow$

\vspace{0.4em}

Helium Coolant

\vspace{0.4em}

$\Downarrow$

\vspace{0.4em}

Steam Generator

\vspace{0.4em}

$\Downarrow$

\vspace{0.4em}

Steam Turbine

\vspace{0.4em}

$\Downarrow$

\vspace{0.4em}

Electricity

\vspace{1em}

\normalsize

We can follow this entire process in the simulator.

\end{frame}

% =============================================================================
\begin{frame}{What Can We Explore in the Simulator?}

\begin{columns}

\begin{column}{0.55\textwidth}

\begin{itemize}
    \item Reactor power
    \item Control rod position
    \item Helium coolant flow
    \item Steam generation
    \item Turbine power output
    \item Electrical generation
    \item System response to changing conditions
\end{itemize}

\vspace{1em}

\begin{alertblock}{Live Demonstration}
Let's change some operating conditions and see how the
plant responds.
\end{alertblock}

\end{column}

\begin{column}{0.45\textwidth}

\centering

% Replace with your controls screenshot
\includegraphics[width=\linewidth]{htr_sim_controls.png}

\end{column}

\end{columns}

\end{frame}

% =============================================================================
\begin{frame}

\centering

\Huge
Live Demonstration

\vspace{2em}

\Large
Questions are welcome throughout.

\end{frame}

% =============================================================================

\end{document}


My actual speaking notes would be:

Slide 1 (60 s)

"This is a high-temperature gas-cooled reactor. It uses helium gas to transport heat."

Slide 2 (90 s)

Point at reactor → helium loop → steam generator → turbine.

Slide 3 (60 s)

"Now let's see what happens when we change reactor power, coolant flow and other operating conditions."

Then immediately switch to the simulator and spend the rest of the station interacting with the schematic. Given the event brief only allocates ~15 minutes per station and asks for a simplified explanation of reactor operation and simulation, this is probably about the right amount of structure.
